\chapter[Sermon 50. The Seventh Angel Blows the Trumpet, and the Elders Sing a Song of Praise.]{The Seventh Angel Blows the Trumpet, and the Elders Sing a Song of Praise.}
\label{sermon:050}
\markright{Sermon 50. The Seventh Angel Blows the Trumpet, and the Elders ...}

The second woe is past, and behold, the third woe will come shortly. And the seventh angel blew, and there were great voices in heaven saying: “The kingdoms of this world have become our lords and Christ’s, and he will reign forevermore.” And the twenty-four elders who sit before God on their thrones fell on their faces and worshiped God, saying: “We give thanks to you, Lord God Almighty, who are, who were, and who are to come; because you have taken hold of your great power and have reigned.”

By the seven angelic trumpets, not only the destinies of the church are shown, but all the faithful are also spurred to watchfulness and to wage spiritual warfare.

\index{Mahomet}\index{Turks}\index{Papacy}\index{Judgment, Last}And to the last three trumpets, as they are the most dangerous, are joined three woes, signifying, as I said at the end of the eighth chapter, that all kinds of troubles and the most grievous afflictions will occur in these times, whereby people will be brought into the greatest distress. And the first, indeed, he has separated from the second and third, by these words: “One woe is past, and lo, two woes are yet to come after this.” This manner of speech does not interrupt the matter, but arranges the speech in order. For the papist woe does not cease when the Turkish woe comes on, but afflicts the churches together. That manner of speaking is given to establish the order; so now he distinguishes the third woe from the second, signifying indeed that Muhammad’s law will endure until the last judgment; and yet in the meantime, he does not deny that papistry will continue also, of which he has hitherto, in the eleventh chapter, discoursed many things, having finished the matters of Muhammad in the ninth chapter. Therefore, the sense of the Apostle’s words seems to be this: “You have heard of the first and second woe, hear furthermore also of the third and last woe.”

And we must note—which greatly contributes to the comfort of the faithful—that the Apostle states plainly that the first and second woes have passed. Thus, he indicates that these two greatest tyrants will have an end, and that God has even appointed them specific limits and boundaries that they cannot exceed. Let us therefore rejoice, that God is watchful over us, and will not abandon us, nor permit the wicked to do more than is fitting.

\index{Daniel (Book of)}The third woe will affect not the godly, but the wicked, at the time they are oppressed by the final judgment and are led astray from their expectations by the Devil to eternal torments. No language, however eloquent, can express the unspeakable pains of this third woe. Therefore, Daniel also says in chapter 12: “The time will be such as has not been since the beginning of humankind.” But when this woe will occur is not expressed nor determined, as neither is the day of judgment, which is known only to the Father and therefore must not be investigated by us with excessive curiosity. It is sufficient for us that it will come shortly. For the Lord says in the gospel that he will shorten those difficult times for the sake of the elect. And again, when these things begin to happen, look up and lift up your heads, for your redemption is drawing near. But these things do not begin now; they are already accomplished. Therefore, it is certain that our redemption is at hand. So cast aside the worry and care with which many torment themselves, thinking that God delays too long, that he allows much to the wicked, that the godly are sorely vexed and, as it were, forsaken and neglected. For the truth says: “Behold, the third woe will come quickly, in due time.” For in chapter 10, he affirmed by a solemn oath that he will come for judgment. Now, as for the precise moment and opportunity of time, give glory to God and acknowledge him in the courses of time and in all things and creatures, using an opportunity most exquisitely. Therefore, when you confess in your creed, “I believe that the Lord will come from the right hand of the Father to judge the living and the dead,” confess also that he will come in due time. And just as from the beginning of the world, he has never forsaken or neglected those who served him, so he will not neglect them at the end of the world.

For it follows that I may expound the things that go before. And the seventh Angel blew. For he declares that the judge is now at hand, raising from the dead the godly and ungodly: the godly unto joy, the wicked to everlasting pain. These shall be new battles, but to the wicked unfortunate, and altogether miserable. Of the trumpet of this Angel, you read in the gospel of Saint Matthew in chapter 24, and in Saint Paul in chapter 4 of the first to the Thessalonians. He should now adjoin the whole manner and discourse of that last judgment, but he will defer it to chapters 19 and 20. In the meantime, he will recite, as he has promised, the furies of Satan against the church, and how he will use those notable instruments, the old and new Roman Empire, to commit murder and, in manner, to destroy the church; wherein, nevertheless, the wicked shall in this world also be put to most grievous punishments. Now omitting, or rather reserving these things to their own place, he celebrates the gratulations, rejoicings, and praises of the Saints.

\index{Augustine, Saint}\index{Rome}The pride and arrogance of the wicked, and especially of the Antichristians, have seemed hitherto intolerable in the world. They have oppressed the godly, boasted of their victories, and bragged of their own felicity. As we shall hear in chapter 18 of this book, that beast has said, “I sit as queen, and am no widow, and shall never see any sorrow.” For voices are heard from Rome: all empires are ours. It is known what manner of things Augustine Steuchus, an Italian and chief champion of the Pope’s authority, has set forth in this cause against Lawrence Valla, about the donation of Constantine. And daily are heard the brags and rejoicings of the papists, concerning the everlasting continuance of the See of Rome, her victories, and the oppression of the preaching of the Gospel, and that she has her power stretched throughout the world, etc.

But in that day, when our Lord Jesus Christ shall abolish all power, rule, and authority, and shall have made all his enemies his footstool, according to the scripture in Psalm 110 and 1 Corinthians 15, there shall be heard again the voices of the glad and joyful, singing true and eternal triumphing songs in heaven. For angels and saints shall sing together; wherefore the voices shall be greater and more durable than the voices of Christ’s enemies, which last but a small season.

Now he also recounts the triumphant song and rejoicing: the kingdoms of this world are made our lords, and Christ’s, and he shall reign forevermore. He shows two things: that all kingdoms are made the fathers’ and the sons’ and that he shall reign forevermore. All kingdoms were before also our Lord Jesus Christ’s; but the same did not appear so plainly to all men, when the bishop of Rome also usurped the same for himself and oppressed those who celebrated only the name of Christ. But when it shall truly appear, and to all flesh, that all kingdoms were ever, and yet remain of one and the eternal God, Christ therefore overcomes, the truth overcomes, the gospel overcomes, the church overcomes: those who are vanquished shall be led to hell, Mahomet with his followers, and the Bishop of Rome with his.

\index{Antichrist}It is added that Christ shall reign forevermore. Antichrist in deed has reigned, and the wicked have rejoiced in this world for a very short time; but now the godly shall reign with Christ forevermore. He does not now divide the kingdom of the Father and the Son, but shows it to be common, where he says that the kingdoms are made: that is to say, it is openly declared that all kingdoms are of God the Father and the Son, and that he shall reign with his elect forevermore. So you may see that the place of Saint Paul may not be explained according to the letter, which is written in 1 Corinthians, chapter 15, concerning the Son being subjected and delivering the kingdom to his Father. For he shall deliver the kingdom, that is to say, shall bring and present it to the Father, and in his members shall be subject to the Father, with whom nevertheless he himself shall reign forever.

The affirmative word “Amen” is annexed, lest any man should doubt one whit of these celestial mysteries. However, he explains afterward more plainly what those voices are that were spoken in Heaven, while he annexes the narration of the twenty-four Elders and of such things wherewith they praised God.

And here the most excellent and beautiful arrangement of this book seems to me worthy of observation. At the beginning of this vision, he introduces the same elders, teaching us by their example and hymns what we should do; therefore, he brings them again also at the end of this vision, that we might be instructed again by their words and actions, not only concerning the final judgment, of what sort it shall be—most righteous, no doubt, as all his judgments are (which the entire vision confirms)—but also that we should understand what is fitting for us, and what we should do: truly, that we should worship God, and wholly submit ourselves to him; and steadfastly believe that both the judgment shall assuredly come, and that it shall also be most just.

The hymn or prayer they offer to God is a form of praise. It is a thanksgiving or rejoicing for victory. In this way, they give God thanks, yet also exalt God and rejoice with themselves and all the faithful in their salvation. They give God thanks for their salvation and commend his justice and truth, which he demonstrates in this judgment, rewarding the good with good things and the evil with evil. Therefore, just as they rise from their chairs and fall down before almighty God, so also we ought to do now and always. This is discussed further in chapter 4. Here we should learn humility, and that God alone is to be worshipped, and that all prayers, invocations, or thanksgivings must be offered to him alone—a matter entirely contrary to the papal doctrine.

We see now the very thanksgiving, for none better can be found. They give thanks unto God. Let us therefore thank him also. And also commend and exalt him while they call him the Lord and God almighty, and also celebrate his majesty, where they say: “Which art, and which wast, and which art to come.” They allude to the words of God, spoken in old time to Moses in Exodus 3. By the diversity of times, the eternity of God is figured. But of this kind of speech I have spoken more in the first Chapter.

\index{Erasmus of Rotterdam}And now they declare why they give thanks: for you have received your great power, and have reigned. God truly never relinquished his power, so that he needs to receive it again; but when he does not display it, and permits much to the ungodly, so that by their power they can infringe and prevail against God’s word, he seems to have laid it aside. Therefore now that he oppresses the wicked, and as a judge advances the godly, maintains the truth, and destroys falsehood, it is truly said that he has received his great power. Likewise now it is said that he reigns, not because he did not reign before, but because the Lord has reigned in the midst of his enemies, so that at times it was doubtful and uncertain whether Christ reigned or Antichrist; yes, that he has had the upper hand, and Christ has been oppressed: now that Christ has broken all the power of his adversaries, it is said most truly that he reigns. And Erasmus rightly admonishes in his annotations on the New Testament, that the translator would have turned more aptly \textit{Ebacilenſas} if he had said, “You have obtained a kingdom.” For the Latin men say, \textit{Regnauit}, “He has reigned,” which has left reigning: as they have lived, who live no more. But with the Greeks it is otherwise, at least in these words. To our judge, most just, most mighty, and most righteous, be praise and glory, forever and ever. Amen.
